% !TeX program = xelatex
\documentclass[UTF8,11pt,a4paper,linespread=1]{ctexart}
\usepackage[margin=26mm]{geometry}
\usepackage{amsmath,amssymb,bm}
\usepackage{setspace}
\setstretch{1.35}
\usepackage[ruled,linesnumbered]{algorithm2e}
\SetKwInput{KwIn}{输入}
\SetKwInput{KwOut}{输出}
\SetArgSty{textnormal}
\SetAlgorithmName{算法}{算法}{算法目录}
\SetAlgoNlRelativeSize{-1}
\SetInd{.5em}{1em}

\begin{document}
\section*{梯度下降：步骤与停止状态分开写}
\begin{algorithm}[htbp]
\SetAlgoLined
\caption{带最大迭代次数的梯度下降}\label{alg:gd}
\KwIn{可微函数 $f$；初值 $\bm{x}_0$；步长 $\eta>0$；
  容差 $\varepsilon>0$；最大迭代次数 $K\in\mathbb{Z}_{\geq 0}$}
\KwOut{最终迭代点 $\bm{x}$、迭代次数 $k$ 与停止原因}
$\bm{x}\gets\bm{x}_0$\;
$k\gets0$\;
$\bm{g}\gets\nabla f(\bm{x})$\;
\While{$k<K$ \textbf{且} $\lVert\bm{g}\rVert_2>\varepsilon$}{
  $\bm{x}\gets\bm{x}-\eta\bm{g}$\;
  $k\gets k+1$\;
  $\bm{g}\gets\nabla f(\bm{x})$\;
}
\eIf{$\lVert\bm{g}\rVert_2\leq\varepsilon$}{
  输出 $\bm{x}$、$k$，停止原因：达到梯度容差\;
}{
  输出 $\bm{x}$、$k$，停止原因：达到迭代上限\;
}
\end{algorithm}
固定步长能否收敛，需要结合函数性质和步长分析。
达到梯度容差也不自动意味着找到了全局最优解。
\end{document}
